\chapter{Conclusion and Future Work}
\label{chap:conclusion}

\section{Conclusion}

This report presented a KGE-NFM pipeline for drug--target interaction
prediction on the Yamanishi08 benchmark. The method first learns fold-specific
drug and protein representations from a biomedical knowledge graph, then
combines these representations with Morgan fingerprints and protein CTD
descriptors in an NFM classifier. DistMult and CompGCN were evaluated as
alternative graph encoders under the same downstream prediction setting.

The experimental results show that adding knowledge graph representations
improves over descriptor-only baselines, and that the relation-aware CompGCN
encoder gives the strongest average performance. CompGCN-NFM obtains
0.9689 ROC-AUC and 0.9681 PR-AUC on Yamanishi 1:1, and 0.9866 ROC-AUC and
0.9362 PR-AUC on Yamanishi 1:10. These values are higher than the corresponding
DistMult-NFM and classical descriptor-based baselines, with the PR-AUC gain on
Yamanishi 1:10 being especially important because that setting contains many
more sampled negatives per positive interaction.

The tuning and ablation results further clarify why the complete framework is
effective. A 200-dimensional CompGCN embedding with five negative samples and
three message-passing layers gave the best tuned result. The ablation study
showed that standalone KGE scoring is not sufficient, while NFM without dense
KGE embeddings is substantially stronger but still below the full model. The
best performance is therefore achieved by combining graph-derived context,
molecular and sequence descriptors, and nonlinear feature interaction modeling.

\section{Key Findings}

\begin{itemize}
    \item The Yamanishi08 DTI matrix is highly sparse, with fewer than one
    percent of all possible drug--target pairs observed as positives.

    \item Knowledge graph embeddings provide useful relational information for
    the downstream DTI classifier.

    \item CompGCN improves over DistMult in the KGE-NFM pipeline, indicating that
    relation-aware neighborhood aggregation is beneficial for this task.

    \item PR-AUC is essential under the Yamanishi 1:10 setting because it exposes
    false-positive control more clearly than ROC-AUC alone.

    \item The ablation results show complementarity: descriptors, KGE embeddings,
    and NFM feature interactions each contribute to the full model.
\end{itemize}

\section{Future Work}

Several directions can strengthen the current study. First, the final report can
incorporate fold-specific confidence intervals or paired significance tests once
the complete fold-level prediction outputs are consolidated. Second, the same
pipeline should be evaluated under drug-cold-start and protein-cold-start
settings to measure generalization to unseen entities. Third, additional KGE
models such as ComplEx, RotatE, and TransE could be compared under the same NFM
prediction protocol.

Representation learning can also be extended beyond the current descriptor
blocks. Molecular graph neural networks may provide richer drug representations
than fixed Morgan fingerprints, while protein language models may capture
sequence patterns not represented by CTD descriptors. Finally, high-confidence
predictions should be validated against external biomedical databases or
literature evidence before being interpreted as biologically actionable
interactions.
